\section{Estructura del documento}
Este documento está dividido en varios capítulos que se centrarán en diversos aspectos del desarrollo, metodología y resultados de este trabajo de fin de estudios. 

En el capítulo \ref{chap-1-Introducción}. Introducción se desarrollan los conceptos relacionados con el preámbulo que justifica el desarrollo de este trabajo. Se abordarán el contexto y motivación del alumno, así como las preguntas de investigación y los objetivos específicos esperados.

En el capítulo \ref{chap-2-Estado-cuestion}. Estado de la cuestión se abordará el conocimiento cientítico existente relacionado con la materia de este TFM. También se abordará los conceptos necesarios para la comprensión del trabajo dentro de su contexto.

En el capiítulo \ref{chap-3-Metodología}. Metodología se recogerán las fases del diseño que da lugar al experimento, la procedencia y características de los datos utilizados y las métricas utilizadas para la evaluación estadística de los resultados.

En el capítulo \ref{chap-4-Diseño-sistema}. Diseño del sistema se abordará en detalle la arquetectura de software donde se han realizado los experimentos.

En el capítulo \ref{chap-5-Experimentos}. Experimentos se recojen cada una de las configuraciones de los modelos desarrollados en la experimentación. Se evaluará el desempeño de cada uno de ellos así como su consistencia.

Análogamente, en el capítulo \ref{chap-6-Resultados-discusion}. Resultados y discusión se presentarán cualitativa y cuantitativamente los resultados obtenidos en los distintos experimentos. También se discutirán las limintaciones de los resultados.

Finalmente, en el capítulo \ref{chap-7-comclusiones}. Conclusiones, se reflexionará acerca de las aportaciones de este trabajo, se responderán las preguntas de investigación planteadas y se introducirán brevemente las posibles líneas de trabajo futuro.

Añadido a esto, se podrá encontrar como secciones independientes, los eanexos y bibliografia del documento.